% page 38 of the facsimile (image p-037 of the scan)
% block 160.0 x 242.0 mm ; left margin 8.0 mm
\pgc{-12.700mm}{0.000mm}{
\l{\cel{3}30}
\l{}
\l{\sou{anodo}. (\sou{elektro}).Elektrodo pozitiva di elektro-baterio,}
\l{\cel{3}di ampulo radiografala. - DEFIRS.}
\l{}
\l{\sou{anomala}.\cel{2}Qua aspektas stranje, baroke. - DEFIRS.}
\l{}
\l{\sou{anonima}. I. Qua ne esas provizita ye la nomo di la autoro.}
\l{\cel{3}- II. Qua ne konocigas sua nomo. - DEFIRS.}
\l{}
\l{\sou{anso}. Parto salianta, ringa od arko, ye qua esas provizita}
\l{\cel{3}utensilo od implemento, por plufaciligar la sizo di olca.}
\l{}
\l{}
\l{\sou{-ant-}\cel{2}Sufixo di la participo aktiva prezenta.}
\l{}
\l{\sou{antagonismo}.\cel{2}Stando di interluktado. - DEFIRS.}
\l{}
\l{\sou{antagonisto}. Ta qua luktas kontre altra. - DEFIRS.}
\l{}
\l{\sou{antarktik}a. Latero qua ne direktesas vers la Urso-stelaro,}
\l{\cel{3}ma vers la sudo-polo di la Terglobo. - DEFIRS.}
\l{}
\l{\sou{ante}.\cel{2}En epoko qua pre-iras la evento quan onu mencionas.}
\l{\cel{3}- DEFIRS.}
\l{}
\l{\sou{antecedento}. I. Antea agi di persono. - II. Termino di pro-}
\l{\cel{3}poziciono qua pre-iras logikale to quo relative \textquotedbl{}konse-}
\l{\cel{3}quas\textquotedbl{}. - III. Termino unesma di relato aritmetikala o}
\l{\cel{3}geometriala. - DEFIS.}
\l{}
\l{\sou{anteno}.\cel{2}I. (\sou{zool.}) Apendico kornatra, artikoza e movebla,}
\l{\cel{3}situita sur la kapo di la maxim multa ek la insekti.}
\l{\cel{3}- II. (\sou{elektro, radio e televiziono}) Longa vergo, fixig-}
\l{\cel{3}ita vertikale relate sulo. - DEFIS.}
\l{}
\l{\sou{antero}. (\sou{bot.}) Mikra posho qua situesas sur la extremajo}
\l{\cel{3}di la stamino e qua kontenas la poleno. - DEFIS.}
\l{}
\l{\sou{anteridio}.\cel{2}(\sou{biol}.)Organo riproduktala maskulala di la}
\l{\cel{3}maxim multa kriptogami, en qua naskas la anterozoidi.-DEFIS.}
\l{}
\l{\sou{anterozoid}o.(\sou{bot}.)Korpi riproduktala maskulala di ula krip-}
\l{\cel{3}togami. - DEFIS.}
\l{}
\l{\sou{antezo}. (\sou{biol}.) Desklozo di la floro, qua liberigas la po-}
\l{\cel{3}leno. - DEFIS.}
\l{}
\l{\sou{anti-}\cel{3}Prefixo (rezervita a la terminaro ciencala e teknik-}
\l{\cel{3}ala) qua signifikas \textquotedbl{}kontre\textquotedbl{}. - DEFIRS.}
\l{}
\l{\sou{anticipar}\cel{2}(\sou{trans.)}\cel{3}I. Prenar, proprigar a su ante la dato.}
\l{\cel{3}- II. Facar, agar, exekutar ante la dato. - DEFIS.}
\l{}
\l{\sou{antidoto}. Substanco quan onu absorbigas da ulu, por kombatar}
\l{\cel{3}en lu la efiko di veneno, di viruso. - DEFIRS.}
}
